% 2007 年四川卷文科第 2 题：四选项图象。
% 来源 PDF：普通高考/2007/2007四川文.pdf，SHA256 1318df44733ec939e8789379d6fa2115a2c6a63e84b73c1f2bf3c080007ce68c
% 原页：第 1 页（A3 三栏）第 1 栏，印刷页 411，图在题干下方。
% 重绘依据：四图 2×2 排列，(A) 递增曲线零点在 x=1、(B) 递减曲线 y 轴截距为 1、(C) 为正确图象（零点 1/2、截距 2、交点 (1,1)）、(D) 两曲线均递增；逐幅与文科原页比对一致。
\begin{tikzpicture}[scale=1.08, semithick, >=Stealth,
  every node/.style={inner sep=1.5pt}]
  % 四个等大的坐标系框（仅平移，画法完全一致）
  \foreach \dx/\dy/\opt in {0/0/A, 4.9/0/B, 0/-4.4/C, 4.9/-4.4/D}{
    \draw[->] ({\dx-0.65},\dy) -- ({\dx+2.85},\dy);
    \draw[->] (\dx,{\dy-0.7}) -- (\dx,{\dy+2.95});
    \node[anchor=north west] at ({\dx+1.07},{\dy-0.07}) {$1$};
    \node[anchor=north west] at ({\dx+2.07},{\dy-0.07}) {$2$};
    \node[anchor=east] at ({\dx-0.28},{\dy+1}) {$1$};
    \node[anchor=east] at ({\dx-0.28},{\dy+2}) {$2$};
    \node[anchor=north east] at ({\dx-0.08},{\dy-0.08}) {$O$};
    \node[anchor=west] at ({\dx+2.92},\dy) {$x$};
    \node[anchor=south west] at (\dx,{\dy+2.97}) {$y$};
    \node[anchor=east] at ({\dx-0.72},{\dy-0.56}) {(\opt)};
  }
  % (A) 递减过 (0,2)、递增过 (1,0)，交于 1<x<2
  \begin{scope}[shift={(0,0)}]
    \draw[domain=-0.35:2.45, samples=100, smooth]
      plot (\x,{2*exp(-0.65*\x)});
    \draw[domain=0.80:2.45, samples=100, smooth]
      plot (\x,{1.55*ln(\x)/ln(2)});
  \end{scope}
  % (B) 递减过 (0,1)、递增过 (0.5,0)，交于 x<1
  \begin{scope}[shift={(4.9,0)}]
    \draw[domain=-0.55:2.60, samples=100, smooth]
      plot (\x,{exp(-0.72*\x)});
    \draw[domain=0.42:2.15, samples=100, smooth]
      plot (\x,{ln(\x/0.5)/ln(2)});
  \end{scope}
  % (C) 真图象：f(x)=1+log_2 x 与 g(x)=2^{1-x}，交于 (1,1)
  \begin{scope}[shift={(0,-4.4)}]
    \draw[domain=-0.30:2.45, samples=100, smooth]
      plot (\x,{exp((1-\x)*ln(2))});
    \draw[domain=0.36:2.45, samples=100, smooth]
      plot (\x,{1+ln(\x)/ln(2)});
  \end{scope}
  % (D) 两条均递增，分别过 (-0.55,0)、(0.42,0)，交于 x≈1.3
  \begin{scope}[shift={(4.9,-4.4)}]
    \draw[domain=-0.64:2.50, samples=120, smooth]
      plot (\x,{1.35*ln(1+(\x+0.55)/1.1)});
    \draw[domain=0.18:1.90, samples=120, smooth]
      plot (\x,{1.60*sign(\x-0.42)*abs(\x-0.42)^1.25});
  \end{scope}
\end{tikzpicture}
